\chapter{Introdução}
\label{cap:introducao}

Organizações dependem de processos documentados para operar de forma previsível, treinar pessoas e cumprir exigências regulatórias. A notação BPMN 2.0, padronizada pela \citeonline{omg2011bpmn}, tornou-se a linguagem de referência para essa documentação, com uso consolidado tanto na indústria quanto na literatura de gestão de processos \cite{weske2019business, dumas2018fundamentals}. Na prática, porém, a maior parte do conhecimento sobre processos permanece registrada em linguagem natural, em manuais internos, descrições de procedimentos e documentos de política, e a conversão desse material em modelos formais continua sendo uma tarefa manual, lenta e dependente de especialistas \cite{friedrich2011process, vanderaa2018challenges}. O problema de extrair estrutura de documentos longos e de vocabulário especializado não é exclusivo da modelagem de processos. No Judiciário brasileiro, por exemplo, a classificação automática de petições iniciais por classe processual tem sido estudada como apoio à distribuição dos processos e ao agendamento de audiências, com modelos de linguagem ajustados ao domínio jurídico \cite{aguiar2021text, aguiar2022topic}.

O avanço dos modelos de linguagem de grande porte \cite{vaswani2017attention, brown2020language} tornou plausível automatizar essa conversão, mas a forma da saída é um obstáculo. A serialização XML do BPMN foi concebida para a troca de arquivos entre ferramentas, e não para ser escrita diretamente por pessoas ou por modelos. Ao pedir a um modelo o XML completo, exige-se dele que reproduza com precisão um grande volume de marcação técnica, embora boa parte dessa marcação possa ser derivada de forma determinística a partir de uma descrição mais curta do processo. Isso tem um custo econômico, porque cada \textit{token} de marcação repetitiva gerado consome tempo e é cobrado. Pode ter também um custo de confiabilidade. Pelas especificações do XML e da linguagem de esquemas, basta uma violação das regras de formação ou do esquema para que o documento inteiro seja rejeitado \cite{bray2008xml, thompson2004xmlschema}, e o modelo gera a saída de forma autorregressiva, um \textit{token} de cada vez \cite{brown2020language}, de modo que cada trecho a mais de marcação é uma oportunidade a mais de erro. O primeiro custo é mensurável diretamente. O segundo é uma hipótese, que os experimentos verificam.

O custo de usar um modelo de linguagem por meio de um provedor é o produto do preço por \textit{token} pelo número de \textit{tokens} processados, e os provedores cobram os \textit{tokens} gerados a um preço várias vezes maior que o dos \textit{tokens} lidos \cite{togetherpricing2026}. Esse custo já pesa nas decisões das empresas. Pode-se citar como exemplo dois casos de 2026. A Uber esgotou em quatro meses o orçamento anual destinado a ferramentas de programação assistida por modelos de linguagem e passou a limitar o gasto a US\$~1.500 mensais por funcionário e por ferramenta, diante da dificuldade, reconhecida pela própria empresa, de relacionar o consumo de \textit{tokens} a um benefício mensurável \cite{ropek2026uber}. No mesmo período, a Microsoft cancelou a maior parte das licenças de assistentes da Anthropic usadas por sua engenharia e passou a direcionar parte das requisições de seus produtos para modelos próprios, com a redução de custo como justificativa declarada \cite{cooter2026microsoft}. Como o custo é o produto do preço pela quantidade, toda forma de reduzi-lo atua sobre um desses dois termos. Algumas atuam sobre o preço efetivo, como o \textit{cache} de trechos repetidos da entrada, o processamento em lote ou a troca de provedor. Outras atuam sobre a quantidade, como encurtar o \textit{prompt}, encurtar a saída ou reduzir as tentativas necessárias até obter uma saída aproveitável. A abordagem investigada aqui atua sobre a quantidade de \textit{tokens} gerados, com o modelo produzindo uma representação mais curta, e sobre o preço, com um modelo de pequeno porte, mais barato, no lugar de um de grande porte.

A redução do número de \textit{tokens} tem sido tratada principalmente do lado da entrada. Na pesquisa, o LLMLingua \cite{jiang2023llmlingua} comprime o contexto antes que ele chegue ao modelo, com ganhos expressivos e pouca perda de desempenho. Na indústria, ferramentas como o Headroom \cite{headroom2026} aplicam a mesma ideia a saídas de ferramentas, registros e trechos recuperados, com reduções relatadas entre 60\% e 95\% em conteúdo estruturado. Essas abordagens mostram que a economia de \textit{tokens} é uma preocupação concreta, mas nenhuma atua sobre a geração, que é a parte mais cara: fazer o próprio modelo gerar uma representação compacta, cuja expansão determinística produza um artefato válido segundo um esquema formal.

Duas linhas de trabalho se aproximam desse objetivo por caminhos diferentes. A decodificação restrita, discutida por \citeonline{poesia2022synchromesh}, garante a conformidade sintática durante a geração, mas não reduz o volume que o modelo precisa gerar. O ProMoAI, de \citeonline{kourani2024promoai}, gera modelos de processo por meio de uma representação intermediária expandida de forma determinística em BPMN, o que mostra que a estratégia é viável. Ele usa, porém, um modelo de grande porte instruído por \textit{prompt} e não mede a compressão que a representação intermediária oferece.

A escolha do modelo também importa. \citeonline{belcak2025small} argumentam que modelos pequenos e especializados podem substituir modelos de grande porte genéricos em tarefas repetitivas e bem delimitadas, por razões de custo, latência e previsibilidade. Uma representação compacta e regular favorece essa especialização, porque reduz o que o modelo precisa aprender a gerar, e torna viável ajustar para a tarefa um modelo de sete bilhões de parâmetros, como o Qwen2.5-Coder \cite{hui2024qwen25coder}.

Propõe-se, assim, uma estratégia de compressão semântica para a geração de BPMN. Um modelo de linguagem descreve o processo em uma linguagem de domínio específico, a BPMN-DSL, e um transpilador determinístico, baseado em uma gramática formal, expande essa descrição no XML do BPMN 2.0 validado contra o esquema oficial. A separação segue o princípio de dividir um sistema em módulos com responsabilidades distintas, que podem ser compreendidos, verificados e alterados de forma independente \cite{parnas1972criteria}. Com isso, uma falha pode ser atribuída à geração da DSL, ao transpilador ou à validação, em vez de ficar diluída em um documento XML extenso. A estratégia é avaliada em processos de um conjunto de teste público, comparando-se modelos de grande porte que geram o XML diretamente, os mesmos modelos gerando a DSL e um modelo pequeno ajustado para a DSL, quanto à validade das saídas, ao custo por requisição e à fidelidade ao texto.

\section{Motivação}
\label{sec:motivacao}

A modelagem de processos ainda é cara justamente onde é mais necessária. Organizações que documentam seus processos em texto raramente dispõem de tempo e de especialistas para convertê-los em modelos formais, e o resultado é uma documentação que não pode ser visualizada, analisada, simulada ou automatizada. Reduzir o custo dessa conversão amplia o acesso à modelagem e libera o especialista para tarefas menos repetitivas.

Há também uma motivação técnica, que não depende do BPMN. Separar a geração de uma representação compacta da sua expansão determinística em um artefato verificável é um padrão aplicável a qualquer domínio em que a saída precise obedecer a um esquema formal e a serialização oficial seja extensa. O BPMN é um caso de teste favorável para essa investigação, porque tem um esquema oficial, que permite verificar a validade de cada saída por programa, e conjuntos de dados públicos com descrições de processos acompanhadas de modelos de referência.

Por fim, a economia de \textit{tokens} deixou de ser um detalhe de implementação. Em fluxos que geram modelos em escala, o custo de gerar marcação redundante se repete a cada requisição, e depender continuamente de modelos de grande porte concentra esse custo em um fornecedor externo. Um modelo pequeno, ajustado para gerar uma representação comprimida, é uma alternativa de engenharia a essa dependência.

\section{Objetivos}
\label{sec:objetivos}

Os objetivos se dividem em um objetivo geral, que formula a pergunta investigada, e em objetivos específicos, que correspondem aos artefatos construídos e às avaliações realizadas para respondê-la.

\subsection{Objetivo Geral}
\label{sec:objetivo-geral}

Investigar se a geração de uma representação intermediária comprimida, seguida de expansão determinística, produz modelos BPMN 2.0 válidos de forma mais confiável e mais econômica em \textit{tokens} do que a geração direta de XML por modelos de linguagem.

\subsection{Objetivos Específicos}
\label{sec:objetivos-especificos}

	\begin{alineas}
		\item projetar uma linguagem de domínio específico para a descrição de processos BPMN, com gramática formalmente especificada;
		\item implementar um transpilador determinístico que expanda essa linguagem no XML do BPMN 2.0, validado contra o esquema oficial;
		\item construir um conjunto de pares de descrição textual e representação na linguagem proposta, a partir de fontes públicas de descrições de processos;
		\item ajustar de forma supervisionada um modelo de linguagem de pequeno porte para gerar a representação comprimida;
		\item comparar a abordagem proposta com a geração direta de XML quanto à validade das saídas, à compressão e ao custo por requisição;
		\item estimar a fidelidade dos modelos gerados ao texto de origem com um modelo de linguagem usado como juiz, submetido a critérios de validação definidos antes de sua aplicação.
	\end{alineas}

\section{Organização do Trabalho}
\label{sec:organizacao}

O Capítulo~\ref{cap:fundamentacao-teorica} apresenta os conceitos usados no restante do texto: a modelagem de processos e a notação BPMN 2.0, os modelos de linguagem e a geração de saídas estruturadas, as linguagens de domínio específico e a transpilação, o ajuste de modelos e os métodos de avaliação e de comparação estatística. O Capítulo~\ref{cap:trabalhos-relacionados} compara a proposta com os trabalhos que geram modelos de processo com modelos de linguagem, diretamente em XML ou por meio de uma representação intermediária, e com as abordagens de compressão de \textit{tokens}. O Capítulo~\ref{cap:metodologia} descreve os dados, a BPMN-DSL, o transpilador, o ajuste do modelo pequeno, os braços experimentais, as medidas e a análise. O Capítulo~\ref{chap:resultados} apresenta a validade, o custo, a compressão e a fidelidade de cada braço, além de uma análise do modelo pequeno especializado. O Capítulo~\ref{chap:conclusoes-e-trabalhos-futuros} responde ao objetivo geral, lista as contribuições e as limitações e propõe trabalhos futuros.
